\section{The extremal setting and the three constraints on \texorpdfstring{\(\nu\)}{nu}}
\label{sec:hypotheses}

We keep the notation of \cite{SS}.  Let \(H=c_{1}(\mathcal
O_{\PP(V)}(1))\), let \(\rho_{1},\dots,\rho_{r}\) and
\(\sigma_{1},\dots,\sigma_{s}\) be the Chern roots of \(V\) and \(V'\)
pulled back to \(\PP(V)\), and let
\(T=\mathrm{tot}(\pi^{*}V'\otimes\mathcal O_{\PP(V)}(-1))\) be the local
model, so that \(c_{1}(T)=(r-s)H+c_{1}(Z)+c_{1}(V)+c_{1}(V')\) by
\cite[(32)]{SS}.

The \(J\)-function of a smooth projective \(Y\) is
\(J^{Y}(Q,t,z)=zS^{Y}(Q,t,z)^{-1}\one\) as in \cite[(20)]{SS}; it satisfies
\begin{equation}
  \label{eq:jnorm}
  J^{Y}(Q,t,z)=z\one+t+O(z^{-1}).
\end{equation}
The \emph{extremal} specialization \cite[(28) and~(33)]{SS} sets \(Q^{d}=1\)
for \(d=k[L]\) with \([L]\) the class of a line in a fibre of
\(\PP(V)\subset T\to Z\), sets \(Q^{d}=0\) for all other \(d\), and replaces
the parameter by
\begin{equation}
  \label{eq:barparam}
  \bar{\mathbf t}=t+\ln(q)c_{1}(T),\qquad t\in H^{*}(Z).
\end{equation}
In the specialized variables \eqref{eq:jnorm} reads
\(J^{T}(q,t,z)=z\one+\bar{\mathbf t}+O(z^{-1})\), and the prefactor
\(z\,e^{t/z}q^{c_{1}(T)/z}=z\,e^{\bar{\mathbf t}/z}\) of \eqref{eq:ss35} is
exactly the \(d=0\) summand.  Whether \eqref{eq:ss35} is the \(J\)-function
is therefore the question of whether the summands with \(d\ge1\) contribute
in \(z\)-order \(\ge0\).

Three places in \cite{SS} carry a hypothesis on \(\nu=r-s\), and the
whole of Sections~5--9 depends on the first and third through the chain
recorded below.  We list each with the inequality it states and the
statements that use it.

\begin{center}
\begin{tabular}{@{}lll@{}}
\hline
Location in \cite{SS} & Hypothesis & Used by\\
\hline
Theorem~4.4 & \(r-s>1\) & (37); Theorem~4.6; (39)--(40)\\
Remark~4.5(3) & \(r-s\le1\) & nothing; a negative assertion\\
Section~7, before (64) & \(r-s>1\) & (64); Proposition~7.5 and its corollaries\\
\hline
\end{tabular}
\end{center}

The first row is the load-bearing one.  \cite[Theorem~4.6]{SS} deduces the
presentation
\begin{equation}
  \label{eq:ss37}
  QH^{*}_{\mathrm{ext}}(T)
  =H^{*}(Z)[H][[q]]\Big/\Big\langle\textstyle\prod_{i=1}^{r}(\rho_{i}+H)
    =q^{r-s}\prod_{j=1}^{s}(\sigma_{j}-H)\Big\rangle,
\end{equation}
which is \cite[(37)]{SS}, from \eqref{eq:ss35} by exhibiting a quantum
differential operator annihilating \(J^{T}\) and quoting
\cite[Theorem~10.3.1]{CoxKatz} and \cite[Lemma~2]{Pandharipande}.  That
argument uses \(r-s>1\) only through \cite[Theorem~4.4]{SS}.  The spectral
computations \cite[Propositions~5.1 and~5.2, Corollary~5.3]{SS} then use
\eqref{eq:ss37} and the ring homomorphism of \cite[Theorem~3.3]{SS}, and
carry no hypothesis on \(\nu\) of their own; the inequality in
\cite[Theorem~1.2]{SS} is inherited from \cite[Theorem~4.4]{SS}.  Likewise,
the modified extremal \(J\)-function \cite[(40)]{SS} used throughout
Sections~7--9 is derived from \eqref{eq:ss35}, and \cite[Proposition~9.1]{SS}
uses \eqref{eq:ss37} again.  Establishing \eqref{eq:ss35} at \(\nu=1\)
therefore unblocks all of this at once.

The second row asserts the opposite of what we prove, and nothing in
\cite{SS} uses it.  The remark then adds that \eqref{eq:ss35} nevertheless
lies on Givental's overruled Lagrangian cone of \(T\) and is therefore an
\(I\)-function; it is this second, positive assertion that the authors
explicitly say they will not need, and the same holds for its counterpart
\cite[Remark~4.7]{SS} for \(T'\).  Cone membership is what we use, and
Section~\ref{sec:normalization} proves it rather than quoting the remark.
The authors attribute it to \cite{Brown}, \cite{CoatesGivental} and
\cite[Lemma~9.6]{SS}; the first two carry no hypothesis on \(\nu\), while
\cite[Lemma~9.6]{SS} is not available for this purpose at \(\nu=1\), for the
reason recorded in Remark~\ref{rem:circularity}.

The third row governs the aperture of the Barnes expansion only.  Nothing in
Sections~8 and~9 of \cite{SS} constrains \(\nu\): the sector
\cite[(78)]{SS} of \cite[Proposition~8.2]{SS} comes from
\cite[Theorem~A.2]{SS}, whose hypothesis \(|\arg t|<(\tfrac{\nu}{2}+1)\pi\)
carries no \(\epsilon\) and is valid at \(\nu=1\); and the reductions of
\cite[Sections~9.1--9.4]{SS} --- weak-to-strong in the local model, the
reduction to the split case, the cohomological reduction to the local model,
and the comparison of the Fourier--Mukai transform with the local model ---
are stated and proved with \(r>s\) as the only inequality.
